\begin{frame}[t]{O Problema da Invertibilidade}
\small

Seja $\KK$ um corpo. Uma \textbf{aplicação polinomial} é uma função
\[
F=(F_1,\ldots,F_n)\colon \KK^n \to \KK^n,
\qquad \text{com } F_i \in \KK[x_1,\ldots,x_n].
\]
Isto é, cada coordenada $F_i$ é um polinômio da forma
\[
F_i=\sum_{\alpha\in \varLambda_i} a_\alpha x^\alpha,
\quad \text{em que } x^\alpha=\prod_{k=1}^n x_k^{\alpha_k},
\]
com $\alpha=(\alpha_1,\ldots,\alpha_n)\in \mathbb{N}^n$, $a_\alpha\in \KK$, e $\varLambda_i\subset \mathbb{N}^n$ um conjunto finito.
\pause
  \begin{defbox}{Pergunta:}
    Quando existe $G\colon\KK^n\to\KK^n$ \textbf{polinomial} tal que
    \(
      F \circ G = G \circ F = \Id_{\KK^n}\,?
    \)
    Qual a relação da inversa com o jacobiano $J_F$?
  \end{defbox}


\end{frame}